\scanpage{98}
% archive foliation: 97
On aura alors
\[
(68)\quad\left\{\begin{aligned}
&\deleted{\unsure{\omega_{i'}}}\quad\omega_{i'}=-\omega_i\\
&\text{i.e. }i\longmapsto\omega_i\text{ est une bijection }\deleted{\unsure{de}}\\
&\alpha=\alpha_\kappa:\mu_4^*\xrightarrow{\sim}\underline\omega
\end{aligned}\right.
\]
% margin position: left, slanted
\margin{\uline{NB} Quand on se donne d'avance
l'iso. $\alpha:\mu_4^*\simeq\underline\omega$, alors
les deux relations (67) sont équiv. \unsure{l'une}\ldots}
On va expliciter la donnée d'un $\kappa$,
\uline{associé} à un isom.
\deleted{$\alpha:\mu_4^*\simeq\underline\omega$}\inserted{donné au préalable}
\[
\alpha=(i\longrightarrow\omega_i):\mu_4^*\simeq\underline\omega,
\]
\deleted{\unsure{$\alpha$}}
via les $\check\kappa_i$. On aura, pour $\delta\in\underline t^*$
\[
(69)\qquad\check\kappa_i(\check\delta)
=\underbrace{\psi_i(\delta)}_{\in F_t}\wedge\omega_i
=\uline{\psi^{\omega_i}(\delta)}\wedge\omega_i,
\]
où
\[
(70)\qquad(\psi^\omega(\delta))_{\omega\in\underline\omega}\in(F_t)^{\underline\omega}
\]
est une famille de \deleted{\unsure{éléments}} deux éléments de $F_t$,
indexée par $\underline\omega$. La condition\inserted{(65) ou} (66)
s'explicite comme
\[
\underbrace{\psi^{\omega_{i'}}(\delta)\wedge\omega_{i'}}_{
\substack{\psi^{\omega'}(\delta)\wedge\omega'\\=-\psi^{\omega'}(\delta)\wedge\omega}}
=\frac1i\underbrace{\psi^{\omega_i}(\delta)\wedge\omega_i}_{\psi^\omega(\delta)\wedge\omega}
\]
\inserted{si \deleted{\unsure{on prend}} $\omega\overset{\mathrm{def}}=\omega_i$, $\omega'=$ autre él. de $\underline\omega$}
% margin position: right
\margin{i.e. $\psi^{\omega'}(\delta)=\underbrace{-\frac1i}_{i}\psi^\omega(\delta)$}
i.e.
\[
(71)\qquad\psi^{\omega'}(\delta)=i_\omega\psi^\omega(\delta)
\qquad\text{où }i_\omega\overset{\mathrm{def}}=\alpha^{-1}(\omega)
\]
On doit avoir de plus
\[
(72)\qquad\psi^\omega(\lambda\delta)=\frac1\lambda\psi^\omega(\delta),
\]
\deleted{\unsure{par ex}} et enfin il faut exprimer (67),
